\chapter{Semilinear Representations and Associated Algebras}\label{chapter:algebras}

Suppose that $G$ is a group, $\bL$ is a field, and $\sigma \colon G \rightarrow \Aut(\bL)$ is group homomorphism from $G$ to the group of field automorphisms of $L$. For $g \in G$ and $\alpha \in \bL$, we often use the shorthand ${}^{g}\alpha\coloneqq \sigma (g) (\alpha)$.
Let $\bK \coloneqq \bL^{\!G}$ be the fixed-point subfield.

The central objects we are interested in are semilinear representations of $G$. In this chapter we review their basic theory.

\section{Semilinear Representations}

\begin{defn}
    A \emph{semilinear representation} of $G$ is a finite dimensional $\bL$-vector spaces $V$ with a group homomorphism $\rho \colon G \rightarrow \Aut_{\bK}(V)$ such that
    \[
    \rho(g)(\alpha \cdot v) = {}^{g}\hspace{-0.1em}\alpha \cdot \rho(g)(v) \ \mbox{  for all  } \ g\in G, \alpha\in\bL , v\in V.
    \]
    A \emph{morphism of semilinear representations} from $V$ to $W$ is an $\bL$-linear map $\phi \colon V \rightarrow W$ such that 
    \[
        \rho_W(g) \circ \phi = \phi \circ \rho_V(g) \ \mbox{  for all  } \ g\in G.
    \]
    We denote the category of such representations by $\Rep_{\bbL}^{\hspace{0.1em} \sigma}(G)$.
\end{defn}

\begin{eg}
    When $\sigma$ is trivial, so $\sigma(g)$ is the identity of $\bL$ for all $g \in G$, $\Rep_{\bbL}^{\hspace{0.1em} \sigma}(G)$ is simply the category of finite-dimensional linear representations of $G$ over $\bL$. In this case we write omit $\sigma$ from the notation and write $\Rep_{\bbL}(G)$ for $\Rep_{\bbL}^{\hspace{0.1em} \sigma}(G)$.
\end{eg}

\begin{defn}
An $\bL$-subspace $W$ of a semilinear representation $V$ is a \emph{subrepresentation} if
\[
        \rho(g)(w) \in W \ \mbox{  for all  } \ g\in G, w \in W.
\]
We say that $V$ is \emph{irreducible} if $V$ is non-zero and $0$ and $V$ are the only subrepresentations of $V$.
We write $\Irr_{\bbL}^{\hspace{0.1em} \sigma}(G)$ for the set of isomorphism classes of irreducible objects of $\Rep_{\bbL}^{\hspace{0.1em} \sigma}(G)$.
\end{defn}

We can perform certain constructions with semilinear representations, which generalise those of linear representations.

\begin{eg}\label{eg:tensorandhoms}
    Suppose that $V, W \in \Rep_{\bbL}^{\hspace{0.1em} \sigma}(G)$. Then
    \begin{itemize}
        \item $V \oplus W$ is a semilinear representation of $G$, where $g \in G$ acts by
    \[
        (v, w) \mapsto (\rho_V(g)(v), \rho_W(g)(w)).
    \]
        \item $V \otimes_{\bbL} W$ is a semilinear representation of $G$, where $g \in G$ acts by
    \[
       v \otimes w \mapsto \rho_V(g)(v) \otimes \rho_W(g)(w).
    \]
    \item $\Hom_{\bbL}(V, W)$ is a semilinear representation of $G$, where $g \in G$ acts by
    \[
        f \mapsto \rho_W(g) \circ f \circ \rho_V(g^{-1}).
    \]
    \end{itemize}
\end{eg}

\begin{eg}\label{eg:trivialrep}
We always have the trivial representation $\bL$, with its natural action of $G$. This is always irreducible and has endomorphism ring $\bK =\bL^{\!G}$.
\end{eg}


\begin{eg}\label{eg:permreps}
Suppose that $G$ acts on a finite set $S$. Then we can form the permutation representation $\bL[S]$, which has $\bL$-basis the elements of $S$, where each $g \in G$ acts via the $\sigma(g)$-linear extension of the permutation action on this basis.
\end{eg}

\begin{eg}\label{eg:inflation}
  Suppose that $f \colon G \twoheadrightarrow Q$ is a surjective morphism of groups, with $\ker(f) \subset \ker(\sigma)$. This kernel condition defines $\sigma$ also on $Q$. There is a pullback or inflation functor
\[
f^* (-) = \Inf_Q^G(-) \colon \Rep_{\bbL}^{\hspace{0.1em} \sigma}(Q) \rightarrow \Rep_{\bbL}^{\hspace{0.1em} \sigma}(G)
\]
defined by letting $G$ act on any $W \in \Rep_{\bbL}^{\hspace{0.1em} \sigma}(Q)$ by the composition of $f$ with $\rho_W \colon Q \rightarrow \Aut_{\bbK}(W)$. 

In particular, this applies to the quotient map $q \colon G \twoheadrightarrow G/N$ for any normal subgroup $N$ of $G$ which is contained in $\ker(\sigma)$. In general, if $V$ is a non-zero semilinear representation of $G$, then $N \coloneqq \ker(\rho_V)$ automatically acts trivially on $\bL$, so $N \subset \ker(\sigma)$, and we can view $V$ as $q^* W$, where $W$ is $V$ with its natural action of $G/N$.
\end{eg}

\section{Matrix Form of Semilinear Representations}\label{sect:matrixform}

We can describe the category $\Rep_{\bbL}^{\hspace{0.1em} \sigma}(G)$ more concretely in terms of matrices. We write $\Vect_{\bbL}$ for the category of finite-dimensional $\bL$-vector spaces.
\subsection{Semilinear Algebra}
For any $\gamma \in \Aut(\bL)$ and $V \in \Vect_{\bbL}$, we write ${}^{\gamma}V$ for $V$ as an abelian group but with $\bL$-vector space structure $\lambda * v = \gamma^{-1}(\lambda)v$. This defines a functor
\[
    {}^{\gamma}(-) \colon \Vect_{\bbL} \rightarrow \Vect_{\bbL},
\] 
which is the identity on morphisms and satisfies ${}^{\gamma}({}^{\mu}V) = {}^{\gamma \mu}V$.

Let $V, W \in \Vect_{\bbL}$ and fix bases $\sA$ and $\sB$ for $V$ and $W$ respectively, which induce an $\bL$-vector space identification
\[
    \End_{\bbL}(V,W) \rightarrow M_{m \times n}(\bL), \qquad [T \colon V \rightarrow W] \mapsto {}_{\sB}[T]_{\sA}.
\]
For any $\gamma \in \Aut(\bL)$, $\sA$ and $\sB$ also form bases for ${}^{\gamma}V$ and ${}^{\gamma}W$ respectively, and there is similarly an $\bL$-vector space identification
\[
    \End_{\bbL}({}^{\gamma}V,{}^{\gamma}W) \rightarrow M_{m \times n}(\bL), \qquad [T \colon V \rightarrow W] \mapsto {}_{\sB}[T]_{\sA},
\]
In fact, there is an equality of sets
\[
    \End_{\bbL}(V,W) = \End_{\bbL}({}^{\gamma}V,{}^{\gamma}W),
\]
and under this equality, the two identifications above are related by the commutativity of the diagram
% https://q.uiver.app/#q=WzAsNCxbMCwwLCJcXEVuZF9SKFYsVykiXSxbMCwxLCJcXEVuZF9SKHt9XntcXGdhbW1hfVYse31ee1xcZ2FtbWF9VykiXSxbMSwxLCJNX3ttIFxcdGltZXMgbn0oUikiXSxbMSwwLCJNX3ttIFxcdGltZXMgbn0oUikiXSxbMywyLCJ7fV57XFxnYW1tYX0oLSkiXSxbMCwzLCJcXHNpbSJdLFsxLDIsIlxcc2ltIl0sWzAsMSwiPSIsMl1d
\[\begin{tikzcd}
	{\End_{\bbL}(V,W)} & {M_{m \times n}(\bL)} \\
	{\End_{\bbL}({}^{\gamma}V,{}^{\gamma}W)} & {M_{m \times n}(\bL)}
	\arrow["\sim", from=1-1, to=1-2]
	\arrow["{=}"', from=1-1, to=2-1]
	\arrow["{{}^{\gamma}(-)}", from=1-2, to=2-2]
	\arrow["\sim", from=2-1, to=2-2]
\end{tikzcd}\]
where ${}^{\gamma}(-) \colon M_{m \times n}(\bL) \rightarrow M_{m \times n}(\bL)$ is the $\bK$-linear automorphism defined by applying $\gamma$ to each entry.

Concretely, if $A \in M_{m \times n}(\bL)$ represents an $\bL$-linear map $T \colon V \rightarrow W$, then ${}^{\gamma \!}A$ represents the $\bL$-linear map ${}^{\gamma}T \colon {}^{\gamma}V \rightarrow {}^{\gamma}W$ (which is set-theoretically simply the map $T$).

Suppose now that $V, W \in \Vect_{\bbL}$, and $S \colon V \rightarrow W$ is a \emph{$\gamma$-semilinear map}, meaning that
\[
S(\alpha \cdot v) = {}^{\gamma \hspace{-0.1em}}\alpha \cdot S(v) \ \mbox{  for all  } \ \alpha \in \bL, v \in V,
\]
or equivalently, that the map $S \colon {}^{\gamma} V \rightarrow W$ is $\bL$-linear. Let $B$ be the matrix associated to the $\bL$-linear map $S \colon {}^{\gamma} V \rightarrow W$ with respect to the bases $\sA$ and $\sB$, and write $[-]$ for the identifications $V \xrightarrow{\sim} \bL^n$, $W \xrightarrow{\sim} \bL^m$ induced by these choices of basis. Then
\[
[S(v)] = B \cdot {}^{\gamma}[v],
\]
where ${}^{\gamma}[v]$ is $\gamma$, applied to all the entries of $[v]$.

\subsection{Semilinear Matrix Representations}
Suppose now that $V \in \Rep_{\bbL}^{\hspace{0.1em} \sigma}(G)$. Then each $\rho(g) \colon V \rightarrow V$ is $\sigma_g$-semilinear and $\rho(g_1 g_2) = \rho(g_1) \circ \rho(g_2)$, so $V$ determines $\bL$-linear maps
\[
(\rho(g) \colon {}^{\sigma(g)}V \rightarrow V)_{g \in G},
\]
which satisfy $\rho(g_1 g_2) =  \rho(g_1) \circ {}^{\sigma(g_1)}\rho(g_2)$. Conversely, given $V \in \Vect_{\bbL}$, any such a family of $\bL$-linear maps defines a semilinear representation structure on $V$.

Let $n \coloneqq \dim_{\bbL}(V)$, and fix a basis of the $\bL$-vector space $V$. For any $g \in G$, write $A_g \in M_{n \times n}(\bL)$ for the matrix of $\rho(g) \colon {}^{\sigma(g)} V \rightarrow V$. The compatibility relations above in matrix form become
\begin{equation}\label{eqn:matrixrelation}
    A_{g_1g_2} = A_{g_1} \cdot {}^{\sigma(g_1) \!} A_{g_2} \ \mbox{  for all  } \ g_1,g_2 \in G.
\end{equation}
In particular, as $I_n = A_{g g^{-1}} = A_{g} \cdot {}^{\sigma(g)} A_{g^{-1}}$, each $A_g \in \GL_n(\bL)$. The following is direct to verify.

\begin{prop}\label{prop:matrixdesc}
Suppose that $V \in \Vect_{\bbL}$ has dimension $n \geq 1$, and fix a basis of $V$.

Then taking the matrix of a semilinear map defines a canonical bijection between semilinear representation structures $\rho$ on $V$ and families $(A_g)_{g \in G}$ in $\GL_n(\bL)$ satisfying (\ref{eqn:matrixrelation}). Under this bijection,
\begin{equation}\label{eqn:matrixevalrelation}
[\rho(g)v] = A_g \cdot {}^{\sigma(g)}[v] \ \mbox{  for all  } \ g \in G, v \in V.
\end{equation}
In particular, an $\bL$-subspace $W$ of $V$ is a subobject of $V$ in $\Rep_{\bbL}^{\hspace{0.1em} \sigma}(G)$ exactly when
\[
        A_g \cdot {}^{\sigma(g)}[w] \in W \ \mbox{  for all  } \ g\in G, w \in W.
\]

Suppose that $W \in \Vect_{\bbL}$ has dimension $m \geq 1$, and fix a basis of $W$. Let $V$ has semilinear representation structure corresponding to $(A_g)_{g \in G}$ in $\GL_n(\bL)$ and let $W$ have semilinear representation structure corresponding to $(B_g)_{g \in G}$ in $\GL_m(\bL)$.

Then an $\bL$-linear morphism $V \rightarrow W$ corresponding to $M \in M_{n \times m}(\bL)$ is a morphism in $\Rep_{\bbL}^{\hspace{0.1em} \sigma}(G)$ if and only if $M$ satisfies
\begin{equation}\label{eqn:matrixhomrelation}
B_g \cdot {}^{\sigma(g)} M = M \cdot A_g \ \mbox{  for all  } \ g \in G.
\end{equation}
\end{prop}

\begin{eg}
Suppose that $V \in \Rep_{\bbL}^{\hspace{0.1em} \sigma}(G)$ has dimension $n \geq 1$, and $\sA$ and $\sB$ are $\bL$-bases of $V$. Suppose that $(A_g)_{g \in G}$ and $(B_g)_{g \in G}$ represent $V$ with respect to $\sA$ and $\sB$ respectively, and that $P \in \GL_n(\bL)$ is the change of basis matrix from $\sA$ to $\sB$, meaning that $P \circ [-]_{\sA} = [-]_{\sB}$. Then
\[
B_g = P \cdot A_g \cdot {}^{\sigma(g)}(P^{-1}).
\]
\end{eg}

\begin{eg}
We can give matrix descriptions of the constructions of Examples \ref{eg:tensorandhoms}, \ref{eg:trivialrep} and \ref{eg:permreps}. The trivial representation has matrix representation by $(I_1)_{g \in G}$, for $I_1$ the identity of $M_1(\bL)$, and if $V$ and $W$ in $\Rep_{\bbL}^{\hspace{0.1em} \sigma}(G)$ have matrix representations $(A_g)_{g \in G}$ and $(B_g)_{g \in G}$ respectively, then 
\begin{itemize}
    \item $V \oplus W$ is represented by the block sum $(A_g \oplus B_g)_{g \in G}$,
    \item $V \otimes_{\bbL} W$ is represented by the Kronecker product $(A_g \otimes B_g)_{g \in G}$,
    \item $\Hom_{\bbL}(V,\bL)$ is represented by $((A_g^{-1})^T)_{g \in G}$.
\end{itemize}

If $G$ acts on a finite set $S$, then $\bL[S]$ has matrix representation by the same matrices as those which represent the linear representation $\bK[S]$, where each $A_g$ is a permutation matrix.
\end{eg}

\section{Semilinear Representations and Skew Group Rings}

Semilinear representations can also be viewed as modules over skew group rings. By $\bL[G]$ we denote the group algebra of $G$.
%%Suppose as above that $\sigma \colon G \rightarrow \Aut(\bL)$ is a group homomorphism. %Write $\bK \coloneqq \bL^{\!G}$, a subfield of $\bL$.

\begin{defn}
    The skew group ring $\bLG$ is $\bL[G]$ as an abelian group, but with multiplication $*$, where
    \begin{equation} \label{LG_product}
        (\alpha_1 \cdot g_1) * (\alpha_2 \cdot g_2) \coloneqq \alpha_1 ({}^{g_1}\alpha_2) \cdot g_1g_2 \ \mbox{  for all  } \ \alpha_1, \alpha_2 \in \bL, \ g_1,g_2 \in G,
    \end{equation}
    extended to all of $\bL[G]$ by $\bZ$-linearly.
\end{defn}

An alternative interpretation of $\bLG$ is that of the set of functions $\varphi : G \rightarrow \bL$ with finite support, where the multiplication formula~\eqref{LG_product} turns into a convolution
    \begin{equation} \label{LG_product2}
      (\varphi \ast \psi) (g) = \sum_{x\in G}  \varphi (x) \; {}^{x}\psi(x^{-1}g) ='' \int_G \varphi (x) \; {}^{x}\psi(x^{-1}g) d\mu_G(x)''.
     \end{equation}
In this section we discuss the properties of $\bLG$ as well as its completions.

\subsection{Categorical properties}
Note that $\bL$ is a subalgebra in $\bL_{\sigma} \hspace{-0.1em}[G]$, but it is not, in general, central.
The central subalgebra is only $\bK$. We therefore naturally view $\bL_{\sigma} \hspace{-0.1em}[G]$ as a $\bK$-algebra.
The following is well-known, and easy to verify.

\begin{prop}
  Suppose that $V \in \Vect_{\bbL}$. Then there is a canonical bijection between left $\bL_{\sigma} \hspace{-0.1em}[G]$-module structures on $V$ and semilinear structures on $V$, under which a morphism in $\Vect_{\bbL}$ is a morphism in $\Rep_{\bbL}^{\hspace{0.1em} \sigma}(G)$ if and only
  if it is a morphism in $\Mod_{\bL_{\sigma} \hspace{-0.1em}[G]}$.
\end{prop}

In particular, we see that $\Rep_{\bbL}^{\hspace{0.1em} \sigma}(G)$ is canonically identified with the full subcategory of $\Mod_{\bL_{\sigma} \hspace{-0.1em}[G]}$ where the underling $\bL$-vector space is finite-dimensional.

\begin{remark}
We note that $\bL$ is still a subalgebra of $\bL[G]$, which is not central whenever the action of $G$ on $\bL$ is non-trivial. Therefore there are two $\bL$-vector space structures on $\bLG$, given by left and right multiplications, and the precise algebraic structure of the skew group ring $\bLG$ is that of a Hopf $\bL|\bK$-algebra \cite{Nich85, Rum00}.
\end{remark}

Let us summarise the tensor properties of semilinear representations in the following, easy to prove, proposition.

\begin{prop} The operations in Example~\ref{eg:tensorandhoms} can be also applied to the category
$\Mod_{\bL_{\sigma} \hspace{-0.1em}[G]}$.
  Under these operations, the following properties hold.
  \begin{enumerate}
  \item $\bL_{\sigma} \hspace{-0.1em}[G]$ is a symmetric closed monoidal $\bK$-linear category.
  \item Its identity object is $\bL$, with the action given by $\sigma$, cf.\ Example~\ref{eg:trivialrep}. 
  \item $\Rep_{\bbL}^{\hspace{0.1em} \sigma}(G)$ is a symmetric  monoidal rigid $\bK$-linear category.
  \item The forgetful functor $\Rep_{\bbL}^{\hspace{0.1em} \sigma}(G)\rightarrow \Vect_{\bbL}$ is tensor $\bK$-linear and faithful, making
    $\Rep_{\bbL}^{\hspace{0.1em} \sigma}(G)$ a Tannakian category \cite{DelMil}.
  \item  The Tannakian category $\Rep_{\bbL}^{\hspace{0.1em} \sigma}(G)$ is neutral if and only if the action $\sigma$ is trivial 
  \end{enumerate}
\end{prop}

We also note here that Maschke's Theorem extends to semilinear representations \cite[Lem.\ 1.3]{KUN}.

\begin{prop}\label{prop:Maschke}
    Suppose $G$ is finite. Then $\Rep_{\bbL}^{\rtimes}(G)$ is semisimple if and only if $|H| \in \bL^{\times}$.
\end{prop}

\subsection{Associated linear representation categories}

There are linear representation categories that interact with $\Rep_{\bbL}^{\hspace{0.1em} \sigma}(G)$. 

\begin{defn}\label{def:H}
We write $H \coloneqq \ker(\sigma \colon G \rightarrow \Aut_{\bbK}(\bL))$.
\end{defn}

With $H$ defined as above, there are $\bK$-algebra inclusions
% https://q.uiver.app/#q=WzAsNSxbMSwxLCJcXGJLW0hdIl0sWzEsMiwiXFxiTFtIXSJdLFsxLDAsIlxcYktbR10iXSxbMiwxLCJcXGJMX3tcXHNpZ21hIFxcaHNwYWNley0wLjFlbX19W0ddIl0sWzAsMSwiXFxiTFtHXSJdLFswLDIsIiIsMCx7InN0eWxlIjp7InRhaWwiOnsibmFtZSI6Imhvb2siLCJzaWRlIjoidG9wIn19fV0sWzAsMSwiIiwyLHsic3R5bGUiOnsidGFpbCI6eyJuYW1lIjoiaG9vayIsInNpZGUiOiJ0b3AifX19XSxbMSwzLCIiLDIseyJzdHlsZSI6eyJ0YWlsIjp7Im5hbWUiOiJob29rIiwic2lkZSI6InRvcCJ9fX1dLFsyLDMsIiIsMCx7InN0eWxlIjp7InRhaWwiOnsibmFtZSI6Imhvb2siLCJzaWRlIjoidG9wIn19fV0sWzIsNCwiIiwwLHsic3R5bGUiOnsidGFpbCI6eyJuYW1lIjoiaG9vayIsInNpZGUiOiJ0b3AifX19XSxbMSw0LCIiLDAseyJzdHlsZSI6eyJ0YWlsIjp7Im5hbWUiOiJob29rIiwic2lkZSI6InRvcCJ9fX1dXQ==
\begin{equation}\label{eqn:diamondrings}\begin{tikzcd}
	& {\bK[G]} \\
	{\bL[G]} & {\bK[H]} & {\bL_{\sigma \hspace{-0.1em}}[G]} \\
	& {\bL[H]}
	\arrow[hook, from=1-2, to=2-1]
	\arrow[hook, from=1-2, to=2-3]
	\arrow[hook, from=2-2, to=1-2]
	\arrow[hook, from=2-2, to=3-2]
	\arrow[hook, from=3-2, to=2-1]
	\arrow[hook, from=3-2, to=2-3]
\end{tikzcd}
\end{equation}
which induce induction and restriction functors between categories
% https://q.uiver.app/#q=WzAsNSxbMSwxLCJcXFJlcF97XFxiYkt9KEgpIl0sWzEsMiwiXFxSZXBfe1xcYmJMfShIKSJdLFsxLDAsIlxcUmVwX3tcXGJiS30oRykiXSxbMiwxLCJcXFJlcF97XFxiYkx9XntcXGhzcGFjZXswLjFlbX1cXHNpZ21hfShHKSJdLFswLDEsIlxcUmVwX3tcXGJiTH0oRykiXSxbMCwyLCIiLDAseyJzdHlsZSI6eyJ0YWlsIjp7Im5hbWUiOiJhcnJvd2hlYWQifX19XSxbMCwxLCIiLDIseyJzdHlsZSI6eyJ0YWlsIjp7Im5hbWUiOiJhcnJvd2hlYWQifX19XSxbMSwzLCIiLDIseyJzdHlsZSI6eyJ0YWlsIjp7Im5hbWUiOiJhcnJvd2hlYWQifX19XSxbMiwzLCIiLDAseyJzdHlsZSI6eyJ0YWlsIjp7Im5hbWUiOiJhcnJvd2hlYWQifX19XSxbMiw0LCIiLDAseyJzdHlsZSI6eyJ0YWlsIjp7Im5hbWUiOiJhcnJvd2hlYWQifX19XSxbMSw0LCIiLDAseyJzdHlsZSI6eyJ0YWlsIjp7Im5hbWUiOiJhcnJvd2hlYWQifX19XV0=
\begin{equation}\label{eqn:diamond}
\begin{tikzcd}
	& {\Rep_{\bbK}(G)} \\
	{\Rep_{\bbL}(G)} & {\Rep_{\bbK}(H)} & {\Rep_{\bbL}^{\hspace{0.1em}\sigma}(G)} \\
	& {\Rep_{\bbL}(H)}
	\arrow[tail reversed, from=1-2, to=2-1]
	\arrow[tail reversed, from=1-2, to=2-3]
	\arrow[tail reversed, from=2-2, to=1-2]
	\arrow[tail reversed, from=2-2, to=3-2]
	\arrow[tail reversed, from=3-2, to=2-1]
	\arrow[tail reversed, from=3-2, to=2-3]
\end{tikzcd}
\end{equation}

For example, in terms of matrices, the induction functor
\[
    \bL_{\sigma \hspace{-0.1em}}[G] \otimes_{\bK[G]} - \colon \Rep_{\bbK}(G) \rightarrow \Rep_{\bbL}^{\hspace{0.1em}\sigma}(G)
\]
simply takes a linear matrix representation $\rho \colon G \rightarrow \GL_n(\bK)$ of $G$ and view this as a semilinear representation under the inclusion $\GL_n(\bK) \hookrightarrow \GL_n(\bL)$. Indeed, the matrices relations (\ref{eqn:matrixrelation}) are satisfied, as all matrices are fixed by $\Aut(\bL / \bK)$.

We will use these linear representation categories to understand $\Rep_{\bbL}^{\hspace{0.1em}\sigma}(G)$.   

\subsection{Bohr completion} \label{Bohr_completion}

In this section, we let $A$ denote a $\bK$-algebra and consider a certain completion of $A$ that allows one to consider simple $A$-modules in terms of idempotents. We will later take $A$ to be one of the algebras of Diagram (\ref{eqn:diamondrings}).

Consider the forgetful functor
\[
\sF : \Mod^{\fd,\semisimple}_A \to \Vect_{\bK}
\]
from the category of finite-dimensional semisimple $A$-modules.
\begin{defn}
  The {\em Bohr completion} of $A$ is $\widehat{A} \coloneqq \End (\sF)$.
\end{defn}

Let $\Irr(A)$ denote the set of isomorphism classes of finite-dimensional simple $A$-modules. For $V \in \Mod^{\fd, \semisimple}_A$, set $E(V) \coloneqq \End_A V$, and let $I(V)$ denote the annihilator of $V$ in $A$.

\begin{lemma} \label{Bohr2}
  Let $\sR$ be a set of representatives of $\Irr(A)$. Then the natural maps
\[ \widehat{A} \xlongrightarrow{\sim}  \prod_{V\in \sR} \End_{E(V)}(V) \xlongleftarrow{\sim} \prod_{V\in \sR} A/I(V)\]
are isomorphisms.
%\marginpar{\textcolor{blue}{I didn't think the surjectivity was obvious, so I added it. We can choose to keep it or not, I personally would. I also added the set $\sR$ to make sure it's clear that this isomorphism is defined for a set of representatives, the product depending on this choice.}}
\end{lemma}

\begin{proof}
  The right-hand map is an isomorphism: the natural map $A \rightarrow \End_{E(V)}(V)$ has kernel $I(V)$, and is surjective when $V$ is irreducible by, for example, the Jacobson Density Theorem.

  For the left-hand map, note that an element of $\widehat{A}$ is a collection $(x_W)$ of $x_W \in \End_{E(W)}(W)$ for each $W \in \Mod^{\fd,\semisimple}_A$, such that $x_{W'} \circ f = f \circ x_{W}$ for any $A$-linear $f \colon W \rightarrow W'$.
  
  If $x_V = 0$ for any $V \in \sR$, then for any $W \in \Mod^{\fd,\semisimple}_A$, $W = W_1 \oplus \cdots \oplus W_n$ for some irreducible submodules $W_i$. Choosing isomorphisms $f_i \colon V_i \xrightarrow{\sim} W_i$ for some $V_i \in \sR$,
  \[
  0 = (f_i \circ \iota_i) \circ x_{V_i} = x_W \circ (f_i \circ \iota_i),
  \]
  where $\iota_i \colon V_i \rightarrow W$ is the inclusion. Then $x_W = 0$, the image of $\oplus_i (f_i \circ \iota_i)$ being all of $W$.
  
  Suppose conversely that $(y_V)_{V \in \sR}$ is any collection of $y_V \in \End_{E(V)}(V)$, subject to no compatibility requirements. To define an extension to a compatible family $(x_W)$ as above, we first set, for $V_1, \ldots , V_n \in \sR$,
  \[
  x_{\underline{V}} \coloneqq x_{V_1} \oplus \cdots \oplus x_{V_n}, \, \text{ where } \, \underline{V} \coloneqq V_1 \oplus \cdots \oplus V_n.
  \]

  Then for any $W \in \Mod^{\fd,\semisimple}_A$ and isomorphism $\phi \colon W \xrightarrow{\sim} \underline{V}$ to such a $\underline{V}$, set
  \[
    x_W^{\phi} \coloneqq \phi^{-1} \circ x_{\underline{V}} \circ \phi.
  \]

 All that remains to show is reduced to the following property: for any morphism $f \colon W \rightarrow X$ in $\Mod^{\fd,\semisimple}_A$, and isomorphisms $\phi \colon W \xrightarrow{\sim} \underline{V}$, $\psi \colon X \xrightarrow{\sim} \underline{U}$, to direct sums of objects of $\sR$, then $f \circ x_X^{\psi} = x_W^{\phi} \circ f$.
  
  Indeed, this shown, then given any $W \in \Mod^{\fd,\semisimple}_A$, we can find some such direct sum $\underline{V}$, and isomorphism $\phi \colon W \xrightarrow{\sim} \underline{V}$, and may set $x_W \coloneqq x_W^{\phi}$. This is independent of the choice of $\underline{V}$ and $\phi$, taking $X = W$ and $f = \id_W$ in the property above, which also tells us that the resulting family $(x_W)_W$ is a well-defined element of $\widehat{A}$.

  To establish the claimed property, note that in the diagram 
  % https://q.uiver.app/#q=WzAsOCxbMCwwLCJcXHVuZGVybGluZXtWfSJdLFszLDAsIlxcdW5kZXJsaW5le1V9Il0sWzAsMSwiXFx1bmRlcmxpbmV7Vn0iXSxbMywxLCJcXHVuZGVybGluZXtVfSJdLFsxLDAsIlciXSxbMSwxLCJXIl0sWzIsMSwiWCJdLFsyLDAsIlgiXSxbNCw3LCJmIl0sWzUsNiwiZiJdLFs2LDMsIlxccHNpIl0sWzcsMSwiXFxwc2kiXSxbMSwzLCJ4X3tcXHVuZGVybGluZXtVfX0iXSxbMCwyLCJ4X3tcXHVuZGVybGluZXtWfX0iXSxbNCwwLCJcXHBoaSIsMl0sWzUsMiwiXFxwaGkiLDJdLFs0LDUsInhfe1d9XntcXHBoaX0iXSxbNyw2LCJ4X3tYfV57XFxwc2l9Il1d
\[\begin{tikzcd}
	{\underline{V}} & W & X & {\underline{U}} \\
	{\underline{V}} & W & X & {\underline{U}}
	\arrow["{x_{\underline{V}}}", from=1-1, to=2-1]
	\arrow["\phi"', from=1-2, to=1-1]
	\arrow["f", from=1-2, to=1-3]
	\arrow["{x_{W}^{\phi}}", from=1-2, to=2-2]
	\arrow["\psi", from=1-3, to=1-4]
	\arrow["{x_{X}^{\psi}}", from=1-3, to=2-3]
	\arrow["{x_{\underline{U}}}", from=1-4, to=2-4]
	\arrow["\phi"', from=2-2, to=2-1]
	\arrow["f", from=2-2, to=2-3]
	\arrow["\psi", from=2-3, to=2-4]
\end{tikzcd}\]
  the commutativity of the central square is equivalent to the commutativity of the outer square. Therefore, it is sufficient to show that for any $A$-linear $g \colon \underline{V} \rightarrow \underline{U}$, $g \circ x_{\underline{V}} = x_{\underline{U}} \circ g$. To see this, label write $\underline{U} = U_1 \oplus \dots \oplus U_m$ for some $U_j \in \sR$. Letting $\iota_i \colon V_i \hookrightarrow \underline{V}$ and $\pi_j \colon \underline{U} \rightarrow U_j$ denote the canonical inclusions and projections, it is sufficient to show that for any $i = 1, ... n$ and $j = 1, ... m$,
  \[
  \pi_j \circ g \circ x_{\underline{V}} \circ \iota_i = \pi_j \circ x_{\underline{U}} \circ g \circ \iota_i.
  \]
  We have that $x_{\underline{V}} \circ \iota_i = x_{V_i} \circ \iota_i$, and $\pi_j \circ x_{\underline{U}} = x_{W_j}\circ \pi_j$, so this is equivalent to the claim that
  \[
  (\pi_j \circ g \circ \iota_i) \circ x_{V_i} = x_{W_j} \circ (\pi_j \circ g \circ \iota_i).
  \]
  If $V_i \neq W_j$, then $V_i \not\cong W_j$ and both maps are zero. If instead $V_i = W_j$, then $\pi_j \circ g \circ \iota_i \in E(V_i)$, and the claim then follows from the fact that $x_{V_i} = x_{W_j} \in \End_{E(V_i)}(V_i)$.
\end{proof}

\begin{remark}
  Lemma~\ref{Bohr2} yields an alternative description of the canonical homomorphism $\iota_A \colon A\to\widehat{A}$, as the map
  \[
  A \rightarrow \prod_{V\in \sR} A/I(V), \qquad a \mapsto (a+I(V))_{V \in \sR}.
  \]
\end{remark}

\begin{eg}
Suppose that $A = \bK[T]$ and $\bK$ is algebraically closed. Then for the theory of the Jordan normal form, there is a canonical identification of $\iota_A \colon A \rightarrow \widehat{A}$ with the product of evaluation homomorphisms
\[
\prod_{\lambda \in \bK} \text{ev}_{\lambda} \colon \bK[T] \hookrightarrow \prod_{\lambda \in \bK} \bK.
\]
\end{eg}

\begin{eg}
The map $\iota_A \colon A \rightarrow \widehat{A}$ doesn't always have to be injective, however. If $A$ is a finite dimensional $\bK$-algebra, then there are finitely many simple modules, and each simple module is finite dimensional, hence $\iota_A \colon A \rightarrow \widehat{A}$ has kernel $J(A)$, the Jacobson radical of $A$, and $\iota_A \colon A \rightarrow \widehat{A}$ is canonically identified with the natural map $A \rightarrow A / J(A)$ to the semisimplification $A / J(A)$ of $A$.
\end{eg}

Not all $\bK$-algebra homomorphisms $\phi : A \rightarrow B$ extend to Bohr completions. This is true, however, in the following setting. Suppose that the restriction functor
\[
(-)|_A \colon \Mod_B \to \Mod_A
\]
preserves finite-dimensional semisimple modules. Then there is a canonical $\bK$-algebra homomorphism $\widehat{\phi} : \widehat{A} \rightarrow \widehat{B}$, defined by
\[
\widehat{\phi}((x_W)_W) = (y_U)_U, \qquad y_U \coloneqq x_{U|_A}.
\]
This is compatible with the canonical homomorphisms $\iota_A \colon A\to\widehat{A}$ and $\iota_B \colon B\to\widehat{B}$ and $\phi \colon A \rightarrow B$.
%\marginpar{\textcolor{blue}{I think this definition is more canonical than the one you gave, as it doesn't need the use the previous lemma and the choice of representatives of the simple modules.}}

We can characterise when $\widehat{\phi} : \widehat{A} \rightarrow \widehat{B}$ is injective as follows.

\begin{lemma} \label{hom_cond}
  Suppose that the functor $(-)|_A \colon \Mod_B \to \Mod_A$ preserves semisimplicity and finite-dimensionality. Suppose also that the induction functor 
  \[ \Mod_A \to \Mod_B, \quad M \mapsto B\otimes_{A}M \]
  preserves semisimplicity and finite-dimensionality. Then $\widehat{\phi} : \widehat{A} \rightarrow \widehat{B}$ is injective if and only if $B \otimes_A V \neq 0$ for any $V \in \Irr(A)$.
\end{lemma}
\begin{proof}
Let $\sR$ be a set of representatives of $\Irr(A)$. Suppose first that the induction functor takes finite-dimensional simple $A$-modules to non-zero (semisimple) $B$-modules, and that $\widehat{\phi}((x_W)_W) = 0$, meaning that $x_{U|_A} = 0$ for any $U \in \Mod^{\fd,\semisimple}_B$. For any $V \in \sR$, there is a canonical $A$-linear map $\iota \colon V \rightarrow B \otimes_A V$ defined by $\iota(v) = 1 \otimes v$. By assumption $B \otimes_A V \neq 0$, and therefore $\iota(V) \neq 0$, hence $\iota$ is injective, by the simplicity of $V$. From the fact that $(x_W)_W \in \widehat{A}$ defines a natural transformation,
\[
\iota \circ x_V = x_{(B \otimes_A V)|_A} \circ \iota = 0,
\]
and therefore $x_V = 0$ from the injectivity of $\iota$. Then $(x_W)_W = 0$, by Lemma \ref{Bohr2}.

Suppose conversely that $\widehat{\phi} : \widehat{A} \rightarrow \widehat{B}$ is injective, and that, for a contradiction, $B \otimes_A V_0 = 0$ for some $V_0 \in \sR$. Consider the element $(x_W)_W \in \widehat{A}$ mapping to $(\delta_{V,V_0})_{V \in \sR}$ under the isomorphism of Lemma \ref{Bohr2}, where $\delta_{V,V} = \id_V$ and $\delta_{V,V_0} = 0$ if $V_0 \neq V$. Then for any $U \in \Mod^{\fd,\semisimple}_B$, 
\[
\Hom_A(V, U|_A) \cong \Hom_B(B \otimes_A V, U) = 0,
\]
and thus for any isomorphism $U|_A \cong V_1 \oplus \cdots \oplus V_n$ with $V_i \in \sR$, each $V_i \neq V$. In particular, as each $x_{V_i} = 0$, then $y_U = x_{U|_A} = 0$, and so, as $U$ was arbitrary, $\widehat{\phi}((x_W)_W) = 0$, whilst $(x_W)_W \neq 0$.
\end{proof}

What is the advantage of $\widehat{A}$ over $A$? It allows us to use idempotents when dealing with simple modules. Let $\sI_A$ denote the set of central primitive idempotents of $\widehat{A}$, and recall that $\Irr(A)$ denote the set of isomorphism classes of irreducible $A$-modules.

Given a simple $A$-module $V$, let $\sR$ be any set of representatives of $\Irr(A)$ containing $V$, and let $e_V \in \widehat{A}$ denote the element of $\widehat{A}$ mapping to $(\delta_{V,V_0})_{V \in \sR}$ under the isomorphism of Lemma \ref{Bohr2}, where $\delta_{V,V} = \id_V$ and $\delta_{V,V_0} = 0$ if $V_0 \neq V$. Concretely, $e_V$ is constructed as in the proof of Lemma \ref{Bohr2}, and acts as the identity on any simple $A$-module $W$ isomorphic to $V$, and as zero on all other simple $A$-modules. In particular, $e_V$ is independent of the choice of representatives $\sR$ and depends only on the isomorphism class of $V$.

\begin{lemma}\label{lem:idemirrcorresp}
  The map $V \mapsto e_V$ defines a bijection $\Irr(A) \rightarrow \sI_A$.
\end{lemma}

\begin{proof}
  The map $V \mapsto e_V$ is injective by construction, and surjective from the description of $\widehat{A}$ of Lemma \ref{Bohr2} as a product of matrix rings over division rings.
\end{proof}

\begin{remark}\label{rem:bohrcompletionexample}
  %\marginpar{\textcolor{blue}{You had written that simple modules for $\widehat{A}$ correspond to idempotents, but I don't think it's quite true. Do you agree with the following?}}
  Given $V \in \Irr(A)$, one can naturally view $V$ as a $\widehat{A}$-module via pullback along the projection $\widehat{A} \twoheadrightarrow \End_{E(V)} (V)$, and in this map we obtain an inclusion
  \[
  \Mod_A^{\fd, \semisimple} / \cong \, \lhook\joinrel\longrightarrow \, \Mod_{\widehat{A}}^{\fd,\semisimple} / \cong.
  \]

  We note that there can be finite-dimensional simple $\widehat{A}$-modules which are not in the image of this map. To construct an example, suppose that $\End_{E(V)}(V) \cong \bF_p$ for any $V \in \Irr(A)$ for some fixed prime $p$, so that any choice $\sR$ of representatives $\sR$ of $\Irr(A)$ induces an isomorphism
  \[
  \widehat{A} \xrightarrow{\sim} \prod_{i \in \sR} \bF_p
  \]
  by Lemma \ref{Bohr2}. Being a commutative ring, the the simple modules of $\widehat{A}$ correspond to the maximal ideals of $\widehat{A}$, which, as $\widehat{A}$ is an product of fields, are naturally identified with ultrafilters on $I$. Under the correspondence the principal ultrafilters correspond the simple modules $V \in \sR$, and the existence of free (i.e.\ non-principal) ultrafilters shows that there exist simple $\widehat{A}$-modules outside the image of the above map whenever $\Irr(A)$ is infinite. To see that these are all finite-dimensional, note that if $\fm$ is any maximal ideal of $\widehat{A}$, then we can view the left $\widehat{A}$-module $M = \widehat{A} / \fm$ as a field extension of $\bF_p$. However, for any surjection 
  \[
  \prod_{i \in \sR} \bF_p \twoheadrightarrow k
  \]
  to a field $k$, $p \cdot k = 0$ and $x^p = x$ for any $x \in k$, hence $k = \bF_p$, and so $M$ is $1$-dimensional over $\bF_p$. 

  To see that such an example is actually possible, consider the group ring $A = \bF_p[G]$, where $G = C_2^{\bN}$ and $p$ is an odd prime. Then finite-dimensional irreducible representations of $G$ over $\bF_p$ are all $1$-dimensional, and correspond to elements of $\{\pm 1\}^{\bN}$, which follows from the fact that an any family of commuting diagonalisable operators on a finite-dimensional vector space can be simultaneously diagonalised. In particular $\Irr(A)$ is infinite, and $\End_{E(V)}(V) = \bF_p$ for all $V \in \Irr(A)$.
\end{remark}

We are interested in $\widehat{A}$ for the following reason. Let us assume that all the maps on Diagram~\eqref{eqn:diamondrings} extend to Bohr completions. Then the completed version, Diagram~\eqref{eqn:diamondrings_complete} below, exhibits five algebras whose central primitive idempotents correspond to the simple finite-dimensional modules of the algebras on Diagram~\eqref{eqn:diamondrings} itself:
\begin{equation}\label{eqn:diamondrings_complete}
 \begin{tikzcd}
	& \widehat{\bK[G]} \\
	\widehat{\bL[G]} & \widehat{\bK[H]} & \widehat{\bL_{\sigma \hspace{-0.1em}}[G]} \\
	& \widehat{\bL[H]}
	\arrow[ from=1-2, to=2-1]
	\arrow[ from=1-2, to=2-3]
	\arrow[ from=2-2, to=1-2]
	\arrow[ from=2-2, to=3-2]
	\arrow[ from=3-2, to=2-1]
	\arrow[ from=3-2, to=2-3]
 \end{tikzcd}
\end{equation}

In Section \ref{sect:blockinterpretation} we will consider central idempotents of $\widehat{\bK[H]}$ which remain central in each larger ring, and use the resulting diagram of algebras, an \emph{S-block}, to study how semilinear representations of $G$ over $\bL$ interact with the linear representation theory of $G$ and $H$ over $\bK$ and $\bL$.

We will use the following lemma, which we advise be skipped and referred back to when needed.

\begin{lemma}\label{lem:gradedzeropartcompletion}
    Suppose that $f \colon A \rightarrow B$ is a $\bK$-algebra morphism, for which the functors
    % https://q.uiver.app/#q=WzAsMixbMCwwLCJcXE1vZF9BIl0sWzEsMCwiXFxNb2RfQiJdLFswLDEsIkIgXFxvdGltZXNfQSAtIiwwLHsiY3VydmUiOi0yfV0sWzEsMCwiKC0pfF9BIiwwLHsiY3VydmUiOi0yfV1d
\[\begin{tikzcd}
	{\Mod_A} & {\Mod_B}
	\arrow["{B \otimes_A -}", curve={height=-12pt}, from=1-1, to=1-2]
	\arrow["{(-)|_A}", curve={height=-12pt}, from=1-2, to=1-1]
\end{tikzcd}\]
    preserve semisimplicity and finite-dimensionality. Suppose that $\phi \colon B \rightarrow B$ is an order two $\bK$-linear automorphism, $\charfield(\bK) \neq 2$, and $f(A) = B^{\phi = 1}$. Then
    \[
    \widehat{f}(\widehat{A}) = (\widehat{B})^{\widehat{\phi} = 1}.
    \]
\end{lemma}

\begin{proof}
    For any $x \in \widehat{A}$ we have that 
    \[
    \widehat{\phi} \widehat{f}(x) = (x_{\phi_* f_* N})_N = (x_{f_* N})_N = \widehat{f}(x),
    \]
    as $f = \phi \circ f$, and thus $\widehat{f}(\widehat{A}) \subset \widehat{B}^{\widehat{\phi} = 1}$. Conversely, suppose that $z \in \widehat{B}$ with $\widehat{\phi}(z) = z$. Because $\phi$ has order two and $\charfield(\bK) \neq 2$, we may define an $A$-linear projection operator
    \[
    \varphi \colon B \rightarrow A, \qquad \varphi(b) = \frac{b + \phi(b)}{2}.
    \]
    We define $x \in \widehat{A}$ by
    \[
        x_M \coloneqq \varphi_M \circ z_{B \otimes_A M} \circ \iota_M,
    \]
    where 
    \[
    \iota_M \colon M \rightarrow B \otimes_A M, \qquad \iota_M(m) = 1 \otimes m
    \]
    is the $A$-linear inclusion and 
    \[
    \varphi_M \colon B \otimes_M \rightarrow M, \qquad \varphi_M(b \otimes m) \coloneqq \varphi(b)m
    \]
    is the $A$-linear map defined using the projection $\varphi$. It is straightforward to see that $x$, defined in this way, gives a well-defined element of $\widehat{A}$, and so it remains to show that $\widehat{f}(x) = z$. This is equivalent to showing that for any $N \in \Mod_B^{\fd,\semisimple}$,
    \[
    x_{f_* N} = \widehat{f}(x)_N = z_N,
    \]
    or equivalently, that
    \[
    \varphi_{f_* N} \circ z_{B \otimes_A N} \circ \iota_{f_* N} = z_N.
    \]
    There is a $B$-linear map
    \[
    \pi_N \colon B \otimes_A N \rightarrow N, \qquad \pi_N(b \otimes n) = bn,
    \]
    and therefore 
    \[
    \pi_N \circ z_{B \otimes_A N} = z_N \circ \pi_N,
    \]
    and pre-composing with $\iota_{f_* N}$ we have that
    \[
    \pi_N \circ z_{B \otimes_A N} \circ \iota_{f_* N} = z_N
    \]
    hence it suffices to show that
    \[
    \varphi_{f_* N} \circ z_{B \otimes_A N} \circ \iota_{f_* N} =  \pi_N \circ z_{B \otimes_A N} \circ \iota_{f_* N}.
    \]
    From the fact that $\widehat{\phi}(z) = z$ means that $z_{\phi_* (B \otimes_A N)} = z_{B \otimes_A N}$. We have a $B$-linear map
    \[
    \phi \otimes 1 \colon B \otimes_A N \rightarrow \phi_*(B \otimes_A N),
    \]
    and therefore (using that $z \in \widehat{B}$),
    \[
    (\phi \otimes 1) \circ z_{\phi_* (B \otimes_A N)} = z_{B \otimes_A N} \circ (\phi \otimes 1)
    \]
    and thus
    \[
    (\phi \otimes 1) \circ z_{B \otimes_A N} = z_{B \otimes_A N} \circ (\phi \otimes 1).
    \]
    By definition, 
    \[
    \varphi_{f_* N} = \frac{1}{2}\left( \pi_N \circ (\phi \otimes 1) + \pi_N \right),
    \]
    and therefore we can compute
    \begin{align*}
       \varphi_{f_* N} \circ z_{B \otimes_A N} \circ \iota_{f_* N} 
       &= \frac{1}{2}\left( \pi_N \circ (\phi \otimes 1) \circ z_{B \otimes_A N} \circ \iota_{f_* N} + \pi_N \circ z_{B \otimes_A N} \circ \iota_{f_* N} \right), \\
       &= \frac{1}{2}\left( \pi_N \circ z_{B \otimes_A N} \circ (\phi \otimes 1) \circ \iota_{f_* N} + \pi_N \circ z_{B \otimes_A N} \circ \iota_{f_* N} \right), \\
    &= \pi_N \circ z_{B \otimes_A N} \circ \iota_{f_* N},
    \end{align*}
    using that $(\phi \otimes 1) \circ \iota_{f_* N} = \iota_{f_* N}$. This shows that $x_{f_* N} = z_N$, and therefore $\widehat{f}(x) = z$, as required.
\end{proof}

\subsection{Bohr Completion for Group Ring Extensions}\label{sect:SblockStructure}
       How common are homomorphisms satisfying the conditions of Lemma~\ref{hom_cond}?
%Recall that an algebra extension $A\subseteq B$ is called {\em separable}. If $B$ is a projective $B\otimes_A B$-module.
      Our goal in this section is to show that the algebra extensions of Diagram~\eqref{eqn:diamondrings} all satisfy these assumptions, subject to mild additional conditions.
       
   \begin{lemma} \label{Sblock_1}
     Suppose $\bK\leq \bL$ is a finite field extension, $A$ is a $\bK$-algebra. Let $B=\bL \otimes_{\bbK} A$ be the extension of scalars.
     Then the restriction functor the restriction functor  $\Mod_B \to \Mod_A$ preserves finite-dimensional semisimple modules.
%     and    defines an algebra homomorphism     $\widehat{A} \rightarrow \widehat{B}$.

If, furthermore, the extension $\bK\leq \bL$ is separable, then the induction functor 
  $\Mod_A \to \Mod_B$
  takes finite-dimensional simple $A$-modules to non-zero finite-dimensional semisimple modules.
%Then the  algebra homomorphism
%    $\widehat{\phi} : \widehat{A} \rightarrow \widehat{B}$ is injective.
\end{lemma}
\begin{proof} 
  Let $V$ be a finite dimension simple $B$-module. The algebra of operators on $V$, coming from $B$ is $B/I(V)$. As in the proof of Lemma \ref{Bohr2}, $B / I(V)$ is isomorphic to a matrix algebra $M_n(D)$, where $D$ is the finite-dimensional division $\bL$-algebra $D \coloneqq \End_B(V)$. Let $A^\prime$ be the image of $A$ in $B/I(V)$. Observe that $A^\prime$ is semisimple: any nilpotent ideal $J$ of $A^\prime$ will give a nilpotent ideal $\bL \otimes_{\bbK} J$ in $B/I(V)$, but $B/I(V)$ is simple, so $J=0$.
  % \marginpar{In fact $A^\prime$ is simple. Do we need it?}
  Thus, $V$ is a semisimple $A$-module.

  For the second statement, consider a finite-dimensional simple $A$-module $M$.
  Similarly, $A/I(M)$ is isomorphic to a matrix algebra $M_n(D)$ for some finite-dimensional division $\bK$-algebra $D$. Under this identification, the induced module $\bL \otimes_{\bK} M$ is identified with $(\bL \otimes_{\bK} D)^n$ and its algebra $B'$ of operators coming from $B$ is identified with $M_n (\bL \otimes_{\bK} D)$.
  Separability implies that $B'$ is semisimple and, hence, the induced $B$-module is also semisimple.
\end{proof}
In light of Lemma~\ref{hom_cond}, we have the following corollary.
\begin{cor}
If $\bK\leq \bL$ is a finite separable field extension, then the canonical $\bK$-algebra homomorphisms
	$\widehat{\bK[H]} \to \widehat{\bL[H]}$ and $\widehat{\bK[G]} \to \widehat{\bL[G]}$ are well-defined and injective.
\end{cor}  

Let $\Gamma$ be a finite group. The second batch of algebras in Diagram~\eqref{eqn:diamondrings} are $\Gamma$-crossed products, for $\Gamma = G/H$.
Recall that if $A$ is a $\bK$-algebra, then a crossed product $A\sharp_{\alpha,\gamma}\Gamma$ is an algebra structure on the group algebra $(A[\Gamma],+)$ given by
\[ ag \cdot bh = a \gamma(g)(b) \alpha(g,h) gh
\ \mbox{ for all } \ a,b\in A, \ g,h \in \Gamma, \]
where $\alpha : \Gamma \times \Gamma \to A$ is a 2-cocycle, $\gamma : \Gamma \to \Aut (A)$ is a homomorphism. This has to satisfy natural conditions ensuring that $A\sharp_{\alpha,\gamma}\Gamma$ is an associative $\bK$-algebra with identity $1_A 1_\Gamma$.

   \begin{lemma} \label{Sblock_2}
     Suppose that $B$ is a crossed product $A\sharp_{\alpha,\gamma}\Gamma$ for some $\bK$-algebra $A$ and finite group $\Gamma$.
      Then the restriction functor  $\Mod_B \to \Mod_A$ preserves finite-dimensional semisimple modules.
%     and    defines an algebra homomorphism     $\widehat{A} \rightarrow \widehat{B}$.

If, furthermore, $\Gamma$ is finite and $|\Gamma| \in \bL^{\times}$, then the induction functor 
  $\Mod_A \to \Mod_B$ takes finite-dimensional simple $A$-modules to non-zero finite-dimensional semisimple modules.
%Then the  algebra homomorphism
%    $\widehat{\phi} : \widehat{A} \rightarrow \widehat{B}$ is injective.
\end{lemma}
   \begin{proof} Observe that the induction of an $A$-module $M$ can be written as
     \[
     B \otimes_A M = \bigoplus_{g\in \Gamma} g \otimes_{\bK} M = \bigoplus_{g\in \Gamma} g \ast M,
     \]
     where the action of $A$ on $g \ast M$ is given by
     \[ a \cdot (g \ast m) = agm = g \gamma(g)^{-1} (a) m = g \ast (\gamma(g)^{-1} (a) \cdot m) . \]
     So it is just the module twisted by the automorphism $\gamma (g)^{-1}$. In particular, if $M$ is semisimple, then $B\otimes_A M$
     is semisimple as an $A$-module.
     
     
     Let $V$ be a finite-dimensional simple $B$-module. There exists a simple $A$-submodule $M$ inside $V$. By simplicity of $V$,
     the induced homomorphism $B\otimes_A M \to V$ is surjective. Thus, $V$ is a quotient of a semisimple $A$-module.
     So it is semisimple as an $A$-module.

     Before proving the second statement, observe that our condition on $\Gamma$ implies that the extension $A \leq B$ is separable in the sense of \cite[Def.\ 2]{HS}. Indeed, its separating idempotent is
     \[ e = \sum_i x_i\otimes_A y_i = \frac{1}{|\Gamma|} \sum_{g \in \Gamma}  \alpha(g,g^{-1})^{-1}g \otimes_A g^{-1} . \]
     This means that any extension of $B$-modules, split as an extension of $A$-modules, is also split as an extension of $B$-modules \cite[Prop.\ 2.6, Thm.\ 1.1]{HS}. Explicitly, given an $A$-splitting $\Phi$, then $\sum_i x_i \Phi y_i$ is the required $B$-splitting (which in our case is just Maschke averaging over the group $\Gamma$).

     Now consider a finite-dimensional simple $A$-module $M$. $B\otimes_A M$ is non-zero by the description above.
     Pick any simple $B$-submodule $V$ inside the induced module $B\otimes_A M$.
  We observed that $B\otimes_A M$ is $A$-semisimple. Thus, the extension $V \hookrightarrow B \otimes_A M$ is $A$-split.
By separability, it is also $B$-split. Hence, $B\otimes_A M$ is $B$-semisimple. 
\end{proof}

\begin{cor}\label{cor:secondbatch}
 If $\bL / \bK$ is finite and $[\bL : \bK] \in \bL^{\times}$, 
 then the canonical $\bK$-algebra homomorphisms
 $\widehat{\bK[H]} \to \widehat{\bK[G]}$,
  $\widehat{\bL[H]} \to \widehat{\bL[G]}$,
 and $\widehat{\bL[H]} \to \widehat{\bL_{\sigma}[G]}$ are well-defined and injective.
\end{cor}  

\begin{proof}
  This follows from Lemma \ref{hom_cond} and Lemma \ref{Sblock_2}, noting that by the assumption that $\bL / \bK$ is finite, $\bL / \bK$ is automatically Galois and the induced map $G / H \rightarrow \Gal(\bL / \bK)$ is an isomorphism \cite[Rem.\ 3.5]{TAYSchur}, hence $\Gamma \coloneqq G/H$ satisfies the assumptions of Lemma \ref{Sblock_2}.
\end{proof}

\begin{remark}
  Actually, for the injectivity of $\widehat{\bL[H]} \to \widehat{\bL_{\sigma}[G]}$, you do not need that $[\bL : \bK] \in \bL^{\times}$, only that $\bL / \bK$ is finite. This will follow from Theorem \ref{thm:facts} below.
\end{remark}

The sixth homomorphism on the Diagram~\eqref{eqn:diamondrings} does not immediately fall into an easily identifiable construction but
we can deal with it by using the result for the other homomorphisms.
% \marginpar{\textcolor{blue}{I think I found a nice proof that doesn't involve taking a normal basis (so doesn't use that $\bL / \bK$ is Galois), but just the result for the other extensions. Does it look ok to you? I put the old proof below if we need to revert back.}}
\begin{thm} \label{SBlock_3}
Suppose that $\bL / \bK$ is finite and $[\bL : \bK] \in \bL^{\times}$.
%Furthermore, assume that $G$ acts on $\bL$ with the kernel $H$ and the extension $\bK\leq \bL$ is $\Gamma$-Galois. (don't think we need this).
Then all homomorphisms in Diagram~\eqref{eqn:diamondrings} satisfy the conditions of Lemma~\ref{hom_cond}, and thus 
all homomorphisms in Diagram~\eqref{eqn:diamondrings_complete} exist and are injective. 
\end{thm}

\begin{proof}
  After the last two lemmas, we know this is true for five of the homomorphisms of Diagram \eqref{eqn:diamondrings},
  and it remains to cobble something together for the sixth homomorphism  ${\bK[G]} \to {\bL_{\sigma}[G]}$. It is clear that induction and restriction between these $\bK$-algebras preserves finite-dimensionality.

  Suppose first that $W$ is a finite-dimensional simple $\bL_{\sigma}[G]$-module. By the commutativity of Diagram \eqref{eqn:diamondrings} and the result for the other homomorphisms of Diagram \eqref{eqn:diamondrings}, $(W|_{\bbK[G]})_{\bbK[H]}$ is semisimple, hence $W|_{\bbK[G]}$ is semisimple being the quotient of the semisimple module
  \[
  \bK[G] \otimes_{\bbK[H]} (W|_{\bbK[G]})_{\bbK[H]} \twoheadrightarrow W|_{\bbK[G]}
  \]
  under the natural surjection.

  If instead $U$ is a finite-dimensional simple $\bK[G]$-module, $\bL_{\sigma}[G] \otimes_{\bK[G]} U$ is non-zero (it can be identified with $\bL \otimes_{\bK} U$ with the diagonal action of $G$ \cite[\S 4.1]{TAYSchur}), and we again have a canonical $\bK[G]$-module surjection
  \[
      \bK[G] \otimes_{\bbK[H]} U|_{\bbK[H]} \twoheadrightarrow U
  \]
  which induces an $\bL_{\sigma}[G]$-module surjection
  \[
  \bL_{\sigma}[G] \otimes_{\bbK[H]} U|_{\bbK[H]} \equiv \bL_{\sigma}[G] \otimes_{\bK[G]} (\bK[G] \otimes_{\bbK[H]} U|_{\bbK[H]}) \twoheadrightarrow \bL_{\sigma}[G] \otimes_{\bK[G]} U
  \]
  from the right-exactness of the tensor product. The $\bL_{\sigma}[G]$-module $\bL_{\sigma}[G] \otimes_{\bbK[H]} U|_{\bbK[H]}$ is semisimple by the commutativity of Diagram \eqref{eqn:diamondrings} and the result for the other homomorphisms of Diagram \eqref{eqn:diamondrings}, and therefore the quotient $\bL_{\sigma}[G] \otimes_{\bK[G]} U$ is too. 
\end{proof}

\section{Semilinear Representations for Finite Galois Extensions}\label{sect:finiteGaloisextensions}

From now on in this chapter we specialise and assume that $\bL / \bK$ is a finite extension. We note that this automatically implies that $\bL / \bK$ is Galois, and that the induced map $\sigma \colon G \rightarrow \Gamma \coloneqq \Gal(\bL / \bK)$ is surjective \cite[Rem.\ 3.5]{TAYSchur}. For example, the assumption that $\bL / \bK$ is finite is automatic if $G$ is finite or $\bL$ is a finite extension of its prime subfield. In this section we recall some results from \cite{TAYSchur} and describe $\Irr_{\bbL}^{\hspace{0.1em} \sigma}(G)$ in terms of $\Irr_{\bbL}(H)$, where $H = \ker(\sigma \colon G \rightarrow \Gal(\bL / \bK))$. 

There is a natural action of $\Gamma$ on $\Irr_{\bbL}(H)$ defined as follows.

\begin{defn}
For $g \in G$ and $V \in \Irr_{\bbL}(H)$ we let $g * V$ be the $\bL[H]$-module which is $V$ as a $\bK$-vector space, but with $\bL[H]$-module structure defined by
\[
\lambda h \cdot v = {}^g\lambda \cdot (g^{-1} h g)(v)
\]
for $\lambda \in \bL$, $h \in H$ and $v \in V$.
\end{defn}
This defines an left action of $G$ on $\Irr_{\bbL}(H)$ which is trivial on $H$, and therefore induces an action of the quotient $G/H$ and thus of $\Gamma$.

\begin{defn}
For $V \in \Irr_{\bbL}(H)$, we let $\Gamma_V$ denote the stabiliser of $V$ in $\Gamma$.
\end{defn}
There are induction and restriction functors between $\Rep_{\bbL}^{\hspace{0.1em} \sigma}(G)$ and $\Rep_{\bbL}(H)$ where
\[
(-)|_H \colon \Rep_{\bbL}^{\hspace{0.1em} \sigma}(G) \rightarrow \Rep_{\bbL}(H)
\]
is just restriction along the inclusion $\bL[H] \rightarrow \bL_{\sigma} \hspace{-0.1em}[G]$ and
\[
\Ind_H^G(-) \coloneqq \bL_{\sigma} \hspace{-0.1em}[G] \otimes_{\bL[H]} - \colon \Rep_{\bbL}(H) \rightarrow \Rep_{\bbL}^{\hspace{0.1em} \sigma}(G).
\]
The following theorem describes how $\Irr_{\bbL}^{\hspace{0.1em} \sigma}(G)$ can be understood in terms of $\Irr_{\bbL}(H)$.
\begin{thm}\label{thm:facts}
Suppose that $\bL/ \bK$ is finite. If $W_1, W_2 \in \Rep_{\bbL}^{\hspace{0.1em} \sigma}(G)$, then:
\begin{itemize}
    \item If $W_1|_H \cong W_2|_H$ then $W_1 \cong W_2$.
    \item The natural map
    \[
    \bL \otimes_{\bbK} \Hom_{\bL_{\sigma} \hspace{-0.1em}[G]}(W_1, W_2) \rightarrow \Hom_{\bL[H]}(W_1|_H, W_2|_H)
    \]
    is an isomorphism.
    \item In particular, if $W \in \Rep_{\bbL}^{\hspace{0.1em} \sigma}(G)$ then the natural map
    \[
        \bL \otimes_{\bbK} \End_{\bL_{\sigma} \hspace{-0.1em}[G]}(W) \rightarrow \End_{\bL[H]}(W|_H)
    \]
    is a $\bK$-algebra isomorphism.
\end{itemize}

Let $W \in \Irr_{\bbL}^{\hspace{0.1em} \sigma}(G)$ and $V \in \Irr_{\bbL}(H)$, where $V$ is an irreducible submodule of $W|_H$, or equivalently that $W$ is an irreducible submodule of $\Ind_H^G(V)$. Then:
\begin{itemize}
    \item $W|_H$ and $\Ind_H^G(V)$ are described by
    \[
        W|_H \cong \bigoplus_{\gamma \in \Gamma / \Gamma_V} (\gamma * V)^{m(W)}, \qquad \Ind_H^G(V) \cong W^{\oplus |\Gamma_V| / m(W)}
    \]
    for some integer $m(W) \geq 1$.
    \item In particular, $W|_H$ and $\Ind_H^G(V)$ are semisimple, and $m(W) \mid |\Gamma_V|$.
    \item $\Ind_H^G(V) \cong \Ind_H^G(\gamma * V)$ for any $\gamma \in \Gamma$.
    \item Any $V \in \Irr_{\bbL}(H)$ is contained in $W|_H$ for a unique $W \in \Irr_{\bbL}^{\hspace{0.1em} \sigma}(G)$. This defines a bijective correspondence
    \[
        \Irr_{\bbL}^{\hspace{0.1em} \sigma}(G) \leftrightarrow \Irr_{\bbL}(H) / \hspace{0.1em}\Gamma.
    \]
\end{itemize}
\end{thm}
\begin{proof}
    This is Theorem A and Theorem B from \cite{TAYSchur}.
\end{proof}

\begin{remark}
  Theorem A and B of \cite{TAYSchur} is stated more generally for indecomposable modules, but we restrict ourselves to the study of irreducible modules here.
\end{remark}

\subsection{Semilinear Schur indices}\label{sect:semilinearschurindex}

Theorem \ref{thm:facts} gives an almost complete description of $\Irr_{\bbL}^{\hspace{0.1em} \sigma}(G)$ purely in terms of $\Irr_{\bbL}(H)$ and its action of $\Gamma$ as follows. 

\begin{defn}
For $V \in \Irr_{\bbL}(H)$, the \emph{semilinear Schur index} of $V$ is $m_{\bbK}^{\bbL}(V) \coloneqq m(W)$, where $W \in \Irr_{\bbL}^{\hspace{0.1em} \sigma}(G)$ is the unique irreducible semilinear $\bL$-representation of $G$ corresponding to $V$.
\end{defn}

Then to construct all objects of $\Irr_{\bbL}^{\hspace{0.1em} \sigma}(G)$, first compute the orbits of the $\Gamma$-action on $\Irr_{\bbL}(H)$. For each such orbit $\sA$, $m_{\bbK}^{\bbL}(V)$ copies of the direct sum
\[
\bigoplus_{V \in \sA} V
\]
extends uniquely to an object of $\Irr_{\bbL}^{\hspace{0.1em} \sigma}(G)$, and all objects of $\Irr_{\bbL}^{\hspace{0.1em} \sigma}(G)$ arise this way.

This description is not quite complete, because it still rests on a determination of the semilinear Schur indices $m_{\bbK}^{\bbL}(V)$ for $V \in \Irr_{\bbL}(H)$, which, like their classical counterparts (of which these are a generalisation \cite[\S 10]{TAYSchur}), appear very hard to compute in general.

It is natural to ask how much information about the numbers $m_{\bbK}^{\bbL}(V)$ is contained in the other associated linear representation theory categories of diagram (\ref{eqn:diamond}). We shall give a satisfactory answer to this question in Chapter \ref{chapter:twelvefoldway}, where we consider the simplest non-trivial case, that $\bL / \bK$ has degree 2.

\subsection{Associated linear representation categories}

The assumption that $\bL / \bK$ is finite allows us to say more about the extensions of diagram (\ref{eqn:diamondrings}). Recall that an extension $R \rightarrow S$ of $\bK$-algebras is called \emph{Frobenius} if the induction and coinduction functors $\Mod_R \rightarrow \Mod_S$ are naturally isomorphic \cite[Def.\ 1.1]{KAD}. The following is Proposition 4.1 of \cite{TAYSchur}.

\begin{prop}\label{prop:frobenius}
If $\bL / \bK$ is finite then all extensions of diagram (\ref{eqn:diamondrings}) are Frobenius.
\end{prop}

\subsection{Induction from subgroups}

For the remained of this chapter we continue with our assumption that $\bL/\bK$ is a finite extension and record some basic properties concerning induction from subgroups. We recommend this be skipped and only referred back to when necessary.

Recall that $\Gamma \coloneqq \Gal(\bL/\bK)$, and $\sigma \colon G \twoheadrightarrow \Gamma$ denotes the action map with kernel $H$. For a subgroup $S$, we can consider the subalgebra $\bL_{\sigma} \hspace{-0.1em}[S]$ of $\bbL_{\sigma} \hspace{-0.1em}[G]$.

\begin{defn}
    For a finite-index subgroup $S \leq G$ we set
    \[
        \Ind_S^G(-) \coloneqq \bL_{\sigma} \hspace{-0.1em}[G] \otimes_{\bbL_{\sigma} \hspace{-0.1em}[S]} (-) \colon \Rep_{\bbL}^{\hspace{0.1em} \sigma}(S) \rightarrow \Rep_{\bbL}^{\hspace{0.1em} \sigma}(G),
    \]
    and for an arbitrary subgroup $S \leq G$ we set
    \[
        (-)|_S \coloneqq \Res_{\bbL_{\sigma} \hspace{-0.1em}[S]}^{\bbL_{\sigma} \hspace{-0.1em}[G]}(-) \colon \Rep_{\bbL}^{\hspace{0.1em} \sigma}(G) \rightarrow \Rep_{\bbL}^{\hspace{0.1em} \sigma}(S).
    \]
\end{defn}

\begin{lemma}\label{lem:inductioncompatibility}
    Let $S$ be a finite-index subgroup of $G$, and $T$ be the kernel of the induced map $S \rightarrow \Gamma$. Then the diagrams
% https://q.uiver.app/#q=WzAsOCxbMCwwLCJcXFJlcF97XFxiYkx9XntcXGhzcGFjZXswLjFlbX0gXFxzaWdtYX0oUykiXSxbMCwxLCJcXFJlcF97XFxiYkx9KFQpIl0sWzEsMCwiXFxSZXBfe1xcYmJMfV57XFxoc3BhY2V7MC4xZW19IFxcc2lnbWF9KEcpIl0sWzEsMSwiXFxSZXBfe1xcYmJMfShIKSJdLFs0LDAsIlxcUmVwX3tcXGJiTH1ee1xcaHNwYWNlezAuMWVtfSBcXHNpZ21hfShHKSJdLFs0LDEsIlxcUmVwX3tcXGJiTH0oSCkiXSxbMywxLCJcXFJlcF97XFxiYkx9KFQpIl0sWzMsMCwiXFxSZXBfe1xcYmJMfV57XFxoc3BhY2V7MC4xZW19IFxcc2lnbWF9KFMpIl0sWzAsMiwiXFxJbmRfe1N9XlQiXSxbMSwzLCJcXEluZF97VH1eSCIsMl0sWzEsMCwiXFxJbmRfe1R9XlMiXSxbMywyLCJcXEluZF97SH1eRyIsMl0sWzQsNywiKC0pfF9TIiwyXSxbNSw2LCIoLSl8X1QiXSxbNyw2LCIoLSl8X1QiLDJdLFs0LDUsIigtKXxfSCJdXQ==
\[\begin{tikzcd}
	{\Rep_{\bbL}^{\hspace{0.1em} \sigma}(S)} & {\Rep_{\bbL}^{\hspace{0.1em} \sigma}(G)} && {\Rep_{\bbL}^{\hspace{0.1em} \sigma}(S)} & {\Rep_{\bbL}^{\hspace{0.1em} \sigma}(G)} \\
	{\Rep_{\bbL}(T)} & {\Rep_{\bbL}(H)} && {\Rep_{\bbL}(T)} & {\Rep_{\bbL}(H)}
	\arrow["{\Ind_{S}^T}", from=1-1, to=1-2]
	\arrow["{(-)|_T}"', from=1-4, to=2-4]
	\arrow["{(-)|_S}"', from=1-5, to=1-4]
	\arrow["{(-)|_H}", from=1-5, to=2-5]
	\arrow["{\Ind_{T}^S}", from=2-1, to=1-1]
	\arrow["{\Ind_{T}^H}"', from=2-1, to=2-2]
	\arrow["{\Ind_{H}^G}"', from=2-2, to=1-2]
	\arrow["{(-)|_T}", from=2-5, to=2-4]
\end{tikzcd}\]
    commute up to natural isomorphism. If additionally $S$ surjects onto $\Gamma$, then the diagrams
% https://q.uiver.app/#q=WzAsOCxbMCwwLCJcXFJlcF97XFxiYkx9XntcXGhzcGFjZXswLjFlbX0gXFxzaWdtYX0oUykiXSxbMCwxLCJcXFJlcF97XFxiYkx9KFQpIl0sWzEsMCwiXFxSZXBfe1xcYmJMfV57XFxoc3BhY2V7MC4xZW19IFxcc2lnbWF9KEcpIl0sWzEsMSwiXFxSZXBfe1xcYmJMfShIKSJdLFs0LDAsIlxcUmVwX3tcXGJiTH1ee1xcaHNwYWNlezAuMWVtfSBcXHNpZ21hfShHKSJdLFs0LDEsIlxcUmVwX3tcXGJiTH0oSCkiXSxbMywxLCJcXFJlcF97XFxiYkx9KFQpIl0sWzMsMCwiXFxSZXBfe1xcYmJMfV57XFxoc3BhY2V7MC4xZW19IFxcc2lnbWF9KFMpIl0sWzAsMSwiKC0pfF9UIl0sWzAsMiwiXFxJbmRfe1N9XlQiXSxbMiwzLCIoLSl8X0giXSxbMSwzLCJcXEluZF97VH1eSCIsMl0sWzQsNywiKC0pfF9TIl0sWzUsNiwiKC0pfF9UIl0sWzYsNywiXFxJbmRfe1R9XntTfSJdLFs1LDQsIlxcSW5kX0heRyIsMl1d
\[\begin{tikzcd}
	{\Rep_{\bbL}^{\hspace{0.1em} \sigma}(S)} & {\Rep_{\bbL}^{\hspace{0.1em} \sigma}(G)} && {\Rep_{\bbL}^{\hspace{0.1em} \sigma}(S)} & {\Rep_{\bbL}^{\hspace{0.1em} \sigma}(G)} \\
	{\Rep_{\bbL}(T)} & {\Rep_{\bbL}(H)} && {\Rep_{\bbL}(T)} & {\Rep_{\bbL}(H)}
	\arrow["{\Ind_{S}^T}", from=1-1, to=1-2]
	\arrow["{(-)|_T}", from=1-1, to=2-1]
	\arrow["{(-)|_H}", from=1-2, to=2-2]
	\arrow["{(-)|_S}", from=1-5, to=1-4]
	\arrow["{\Ind_{T}^H}"', from=2-1, to=2-2]
	\arrow["{\Ind_{T}^{S}}", from=2-4, to=1-4]
	\arrow["{\Ind_H^G}"', from=2-5, to=1-5]
	\arrow["{(-)|_T}", from=2-5, to=2-4]
\end{tikzcd}\]
    commute up to natural isomorphism.
\end{lemma}

\begin{proof}
    The commutativity of the first two diagrams follows directly from the transitivity of induction and restriction. Suppose now that the induced map $\sigma \colon S \rightarrow \Gamma$ is surjective, which means that $[G:H] = [S:T]$ and $[G:S] = [H:T]$. The first isomorphism is given for $W \in \Rep_{\bbL}^{\hspace{0.1em} \sigma}(S)$ by
    \[
        \phi \colon \bL[H] \otimes_{\bbL[T]} (W|_T) \rightarrow \left( \bL_{\sigma} \hspace{-0.1em}[G] \otimes_{\bbL_{\sigma} \hspace{-0.1em}[S]} W \right)|_H, \qquad \phi(h \otimes w) \coloneqq h \otimes w.
    \]  
    Because $[G:S] = [H:T]$, both sides have $\bL$-dimension $[G:S] \dim_{\bbL}(W)$, and so it is sufficient to show that $\phi$ is surjective. Indeed, for any $g \in G$, $g = h s$ for $h \in H$ and $s$ in $S$, because the coset representatives of $T$ in $S$ are coset representatives for $H$ in $G$, and so for any $w \in W$,
    \[
        \phi(h \otimes s(w)) = h \otimes s(w) = g \otimes w.
    \]
    The second isomorphism is similarly given for $V \in \Rep_{\bbL}(H)$ by
    \[
        \psi \colon \bL_{\sigma} \hspace{-0.1em}[S] \otimes_{\bbL[T]} (V|_T) \rightarrow \left( \bL_{\sigma} \hspace{-0.1em}[G] \otimes_{\bbL[H]} V \right)|_S, \qquad \phi(s \otimes v) \coloneqq s \otimes v,
    \]  
    which can similarly be seen to be an isomorphism using that any $g \in G$ has the form $g = s h$ for $s \in S$ and $h \in H$ and that both sides have the same $\bL$-dimension as $[G:H] = [S:T]$.
\end{proof}

\begin{lemma}\label{lem:projectionformula}
    Suppose that $S \leq G$ is a finite-index subgroup. Then there is a natural isomorphism
    \[
        \Ind_S^G(W \otimes_{\bbL} V|_S) \xrightarrow{\sim} \Ind_S^G(W) \otimes_{\bbL} V
    \]
    in $V \in \Rep_{\bbL}^{\hspace{0.1em} \sigma}(G)$ and $W \in \Rep_{\bbL}^{\hspace{0.1em} \sigma}(S)$.
\end{lemma}

\begin{proof}
    The natural map of $\bL_{\sigma} \hspace{-0.1em}[S]$-modules
    \[
        W \otimes_{\bbL} V|_S \rightarrow (\bL_{\sigma} \hspace{-0.1em}[G] \otimes_{ \bbL_{\sigma} \hspace{-0.1em}[S]} W) \otimes_{\bbL} V,
    \]
    extends to
    \[
    \bL_{\sigma} \hspace{-0.1em}[G] \otimes_{ \bbL_{\sigma} \hspace{-0.1em}[S]}(W \otimes_{\bbL} V|_S) \rightarrow (\bL_{\sigma} \hspace{-0.1em}[G] \otimes_{ \bbL_{\sigma} \hspace{-0.1em}[S]} W) \otimes_{\bbL} V
    \]
    which is surjective, hence an isomorphism, both sides having the same $\bK$-dimension.
\end{proof}

\subsubsection{Comparison with Classical Induction}

We also record here that this induction is compatible with induction in the split case, under the equivalence of \cite[\S 6]{TAYSchur}.

\begin{lemma}\label{lem:compatiblewithsplitinduction}
    Suppose that $S$ is a subgroup of $H$. Then the diagrams
% https://q.uiver.app/#q=WzAsOCxbMCwwLCJcXFJlcF97XFxiYkt9KEgpIl0sWzAsMSwiXFxSZXBfe1xcYmJLfShTKSJdLFsxLDAsIlxcUmVwX3tcXGJiTH1ee1xcaHNwYWNlezAuMWVtfSBcXHNpZ21hfShIIFxcdGltZXMgXFxHYW1tYSkiXSxbMSwxLCJcXFJlcF97XFxiYkx9XntcXGhzcGFjZXswLjFlbX0gXFxzaWdtYX0oUyBcXHRpbWVzIFxcR2FtbWEpIl0sWzMsMCwiXFxSZXBfe1xcYmJLfShIKSJdLFs0LDAsIlxcUmVwX3tcXGJiTH1ee1xcaHNwYWNlezAuMWVtfSBcXHNpZ21hfShIIFxcdGltZXMgXFxHYW1tYSkiXSxbMywxLCJcXFJlcF97XFxiYkt9KFMpIl0sWzQsMSwiXFxSZXBfe1xcYmJMfV57XFxoc3BhY2V7MC4xZW19IFxcc2lnbWF9KFMgXFx0aW1lcyBcXEdhbW1hKSJdLFswLDIsIlxcc2ltIl0sWzEsMywiXFxzaW0iXSxbNiw3LCJcXHNpbSJdLFs0LDUsIlxcc2ltIl0sWzEsMCwiXFxJbmRfe1N9XkgiXSxbMywyLCJcXEluZF97UyBcXHRpbWVzIFxcR2FtbWF9XntIIFxcdGltZXMgXFxHYW1tYX0iLDJdLFs0LDYsIigtKXxfUyIsMl0sWzUsNywiKC0pfF97UyBcXHRpbWVzIFxcR2FtbWF9Il1d
\[\begin{tikzcd}
	{\Rep_{\bbK}(H)} & {\Rep_{\bbL}^{\hspace{0.1em} \sigma}(H \times \Gamma)} && {\Rep_{\bbK}(H)} & {\Rep_{\bbL}^{\hspace{0.1em} \sigma}(H \times \Gamma)} \\
	{\Rep_{\bbK}(S)} & {\Rep_{\bbL}^{\hspace{0.1em} \sigma}(S \times \Gamma)} && {\Rep_{\bbK}(S)} & {\Rep_{\bbL}^{\hspace{0.1em} \sigma}(S \times \Gamma)}
	\arrow["\sim", from=1-1, to=1-2]
	\arrow["\sim", from=1-4, to=1-5]
	\arrow["{(-)|_S}"', from=1-4, to=2-4]
	\arrow["{(-)|_{S \times \Gamma}}", from=1-5, to=2-5]
	\arrow["{\Ind_{S}^H}", from=2-1, to=1-1]
	\arrow["\sim", from=2-1, to=2-2]
	\arrow["{\Ind_{S \times \Gamma}^{H \times \Gamma}}"', from=2-2, to=1-2]
	\arrow["\sim", from=2-4, to=2-5]
\end{tikzcd}\]
commute up to natural isomorphism, where the horizontal arrows are the equivalences of \cite[\S 6]{TAYSchur}.
\end{lemma}

\begin{proof}
    For $U \in \Rep_{\bbK}(S)$, the first isomorphism is given by
    \[
        \bL \otimes_{\bbK} (\bK[H] \otimes_{\bbK[S]} U) \rightarrow \bbL_{\sigma} \hspace{-0.1em}[H \times \Gamma] \otimes_{\bbL_{\sigma} \hspace{-0.1em}[S \times \Gamma]} (\bL \otimes_{\bbK} U), \qquad \lambda \otimes h \otimes u \mapsto h \otimes \lambda \otimes u
    \]
    for $\lambda \in \bL$, $h \in H$ and $u \in U$. This is clearly surjective, hence an isomorphism as both sides have the same $\bL$-dimension. The second isomorphism is clear.
\end{proof}
